\documentclass{article}

\usepackage{booktabs}
\usepackage{tabularx}

\title{Development Plan\\\progname}

\author{\authname}

\date{}

\input{../Comments.text}
\input{../Common.text}

\begin{document}

\maketitle

\begin{table}[hp]
\caption{Revision History} \label{TblRevisionHistory}
\begin{tabularx}{\textwidth}{llX}
\toprule
\textbf{Date} & \textbf{Developer(s)} & \textbf{Change}\\
\midrule
Date1 & Name(s) & Description of changes\\
Date2 & Name(s) & Description of changes\\
... & ... & ...\\
\bottomrule
\end{tabularx}
\end{table}

\newpage{}

\section{Introduction}

This document outlines the proposed development plan for the team. Sections 1-3 contain legal and regulatory information, sections 4-6 contain logistical details for team communication and roles, and Section 7-10 contain timeline
and technology logistics. The team charter is also included at the end of this document.

\section{Confidential Information?}

The images used to train the machine learning model are confidential and cannot be shared publicly. The images will be provided by various municipalities and are subject to a confidentiality agreement.
The team will not share any images or data that is considered confidential. A confidentiality agreement has not been completed yet, but a template is available via this link: \href{../../Capstone Template - Confidentiality Only.docx}{Capstone Template}.

\section{IP to Protect}

For this project, there is IP to protect. An IP agreement has not been completed yet, but a template is available via this link: \href{../../Capstone Template - IP Owned By Company.docx}{Capstone Template}.

\section{Copyright License}

The team is explicity choosing not to include a copyright license for this project, meaning that the code is under exclusive copyright. This information is also available in the license file: \href{../../LICENSE}{LICENSE}.

\section{Team Meeting Plan}

Team members will meet on a weekly basis on Tuesdays from 11 AM - 12:30 PM. Meetings will
be located in-person at one of McMaster's libraries. Team members are expected to come
prepared by having reviewed the current state of the deliverable at that time. Meetings may
occur more than once per week if needed, which will be held virtually on Discord. Supervisor meetings
will be held regularly on a biweekly basis, but the team will be open to meeting more frequently if
the supervisor requests it. All members will take turns chairing meetings and taking meeting notes. 
Meetings should be structured as follows.
\newline

At the start of every meeting, the chair will go over the purpose of the meeting and discuss what is on the agenda. 
The meeting notetaker should take note of the agenda.
\newline

After outlining the agenda, the chair will go through each meeting item one-by-one. Each item should 
have a person responsible for that task to report on progress, expected outcomes, and 
explain next steps if necessary. Each item should have a clear time allocation, and expected
deadline. The meeting notetaker should take note of all of this information.
\newline

At the end of the meeting, the chair will summarize all key decisions made, and restates
actions items for the week. The meeting notetaker will write a section for meeting summary notes.

\section{Team Communication Plan}

Team will use the following communication methods:

\begin{itemize}
    \item Issues: GitHub Issues
    \item Regular Weekly Meetings: In-person at McMaster Libraries
    \item Virtual Meetings: Discord
    \item Stakeholder Meetings: Teams
    \item Meeting minutes \& Agenda: GitHub
    \item Project discussion: Discord
    \item Urgent/Emergency Communication: SMS Messaging
\end{itemize}

Team members will regularly check in on communication channels for updates,
and enable notifications. Team members should make an effort to respond to messages within
24 hours. Team members will exchange phone numbers to communicate in case of emergency or urgent situations.

\section{Team Member Roles}

\begin{itemize}
    \item{Note Taker: Takes notes to document key events that occurred during weekly team meetings,
 summarizing key decisions made during the meeting. Responsible for producing 
organized, neat, and concise notes made accessible to the team for reference.\newline}
    \item {Meeting Chair: Is responsible for leading discussion for weekly team meetings. Should outline a
meeting agenda in advance to facilitate smooth discussion of current issues and action items. \newline}
    \item {Task Owner: Responsible for complete ownership of a particular task or deliverable. During a team meeting,
they are expected to know important details regarding a task's progress and expected deadlines. They should 
be prepared to report on a task during weekly stand-up.}
    \item{LLaison: Responsible for external communication with faculty and all 3rd parties.\newline}
\end{itemize}

\section{Workflow Plan}

Git will be used for version control and collaboration. The team will use GitHub to host the repository and manage issues, pull requests, and code reviews. Each team member will be responsible for creating their own branches for their work and submitting pull requests for review before merging into the main branch.
A general workflow plan is listed below.

\begin{itemize}
  \item When a deliverable or task is necessary, an issue will be created in the GitHub repository. The issue will contain a description of the task, expected outcomes, and any relevant information. The issue will be assigned to a team member who will be responsible for completing the task.
  \item The team member will create a new branch for the task and work on it. The branch should be named according to the issue number and a brief description of the task.
  \item Once the task is completed, the team member will submit a pull request for review. The pull request should contain a description of the changes made, any relevant information, and a link to the issue.
  \item The team member will request a review from at least one other team member. The reviewer will review the code, provide feedback, and approve or request changes to the pull request.
  \item Once the pull request is approved, the team member will merge the changes into the main branch. The issue will be closed, and the task will be considered complete.
\end{itemize}

GitHub Issues will be used to track deliverables/tasks, lectures, and supervisor/TA/team meetings. For deliverables and tasks, tags will be used to identify the kind of task that needs to be done. This could include tags like "bug fixing", "feature development", "documentation", etc. For meetings, tags will be used to identify the type of meeting.

GitHub Actions will be used as the CI/CD platform, and will be used for linting, unit testing, and compiling any documentation written in LaTeX. These workflows will be automatically on push and pull requests, and ensure that both code and documentation are up to standard.


\section{Project Decomposition and Scheduling}

The team will use Github Projects to track tasks and deliverables. The team will use a Kanban board to track the progress of tasks and deliverables. The board will have columns for "To Do", "In Progress", "In Review", and "Done".
Each task or deliverable will be represented by a card on the board, which will contain a description of the task, expected outcomes, and any relevant information. The cards will be moved across the columns as the task progresses.
Each card will be tied to a GitHub issue, which will be used to track the progress of the task. The issue will be assigned to a team member who will be responsible for completing the task.

The project schedule below follows the timeline detailed in the course syllabus.
 \begin{itemize}
   \item Weeks 4: Problem Statement, Proof of Concept (POC) Plan, Development Plan, Team Charter
   \item Week 6: Requirements Document (SRS) Revision 0, Hazard Analysis Revision 0
   \item Week 8: Verification \& Validation (V\&V) Plan Revision 0
   \item Week 10: Design Document Rev -1 and CD-Fall Rev 0
   \item Weeks 11+12: Proof of Concept Demo + Interviews
   \item Week 16: Design Document Revision 0
   \item Week 18+19: Revision 0 Demo + Interviews
   \item Week 22: V\&V Report Rev 0 and CD-Winter Rev 0
   \item Week 24: Final Demonstration (Revision 0)
   \item Week 26: EXPO Demonstration, Final Documentation (Revision 1)
 \end{itemize}

\section{Proof of Concept Demonstration Plan}

The main risk towards the success of the project is the ability to accurately detect and classify the state of the road based on the information given.
During the proof of concept demonstration, we will demonstrate that our model can be given an image of a road and accurately detect and classify the state of the road.

Some major risks and corresponding mitigation strategies are as follows:
\begin{itemize}
    \item Risks associated with the quality and quantity of images provided by municipalities. Currently, the project assumes that the team will have a large, labelled dataset of images to train the model with.
    \begin{itemize}
      \item To address this, the team will do research into public datasets that can be used to train the model. The team will also investigate data augmentation techniques to increase the size of the dataset.
    \end{itemize}  
    \item Another major risk is the model's ability to classify the state of the road in a way that is universilly applicable to all municipalities. The model may be able to classify the state of the road in a way that is specific to a certain municipality, but not applicable to others.
    \begin{itemize}
      \item As a midway point, the model will be trained to classify the state of the road in a way that is applicable to a specific municipality. The team will then investigate how to generalize the model to work for other municipalities.
    \end{itemize}
\end{itemize}

\section{Expected Technology}

Please note that the technologies and tools listed above are not final and may change as the project progresses. The team will evaluate the effectiveness of the tools and technologies used and make adjustments as necessary.

Python will be the primary programming language used for this project. We will be using libraries such as OpenCV for image processing, TensorFlow or PyTorch for machine learning, and NumPy for numerical computations. We will also utilize pre-trained models for road state detection and classification.

Ruff will be used as the linter for Python code, and PyTest will be used as the unit testing framework. We will investigate code coverage measuring tools such as Coverage.py to ensure that our tests are comprehensive.

For version control, we will use Git and GitHub. We will use CI for linting and unit test passing. Since we are using a Github repo, Github Actions is the preferred method of choice for CI. 

\section{Use of AI and LLMs}

Ai and LLMs will be used to assist in the development of code and debugging. The team will also use AI to assist in the writing of documentation. The agents that the team will use are Github Copilot, Claude Sonnet 5.5, and Google Gemini Flash 3.6.
All code and documentation generated by AI will be reviewed by at least two members of the team to ensure accuracy and quality, and must pass linting and unit tests before being pushed into a branch. Additionally, planning mode will be used for any
agents that are used to ensure that the output is accurate and relevant to the task at hand.

\section{Coding Standard}

PEP 8 will be used as the coding standard for Python code. In the event that a different language is used, the team will abide by the standards set by the creators of the language.

In the event that the coding standard fails to cover a specific case, the team will discuss and agree on a standard to follow for that specific case. Linters will be used to ensure consistent and clean code according to the specified standards.

\newpage{}

\section{Appendix --- Generative AI Use Disclosure}

\subsection{AI Tools Used}

Github Copilot, version 1.0.88.

\subsection{How and Where Used}

Github Copilot was used to complete sentences in the use of AI + LLM, expected technology, and workflow plan sections.

\subsection{Verification of AI-Generated Content}

Since Copilot only completes sentences, the team wrote short, concise sentences that were then expanded by Copilot. This way, any ideas in the document are the team's own, and Copilot is only used to expand on those ideas.
For any documentation generated by Copilot, a second member of the team reviewed the content before allowing it to be pushed. For any technologies that were suggested by Copilot, the team manually researched the technology to ensure that it was relevant and applicable to the project.

\newpage{}

\section*{Appendix --- Reflection}

\wss{Not required for CAS 741}

\input{../Reflection.text}

\begin{enumerate}
    \item Why is it important to create a development plan prior to starting the
    project?
    \item In your opinion, what are the advantages and disadvantages of using
    CI/CD?
    \item What disagreements did your group have in this deliverable, if any,
    and how did you resolve them?
\end{enumerate}

\newpage{}

\section*{Appendix --- Team Charter}

\subsection*{External Goals}

\wss{What are your team's external goals for this project? These are not the
goals related to the functionality or quality fo the project.  These are the
goals on what the team wishes to achieve with the project.  Potential goals are
to win a prize at the Capstone EXPO, or to have something to talk about in
interviews, or to get an A+, etc.}

Our team's external goals for this project are to gain experience in the space of machine learning and work on a project with real-world applications and stakeholders.
Additionally, we hope to develop a project that can be a strong addition to our portfolios.

\subsection*{Attendance}

\subsubsection*{Expectations}

For planned team meetings, all members are expected to attend either online or virtually. If a team member is unable to attend, they must notify the team at least 12 hours in advance.
The team will wait up to 5 minutes for late members to arrive, and will start the meeting after this time period. After arriving, team members are expected to stay for the duration of the meeting,
and should only leave early for exceptional circumstances.

\subsubsection*{Acceptable Excuse}

Our team will operate on the principle that we are all responsible for our own time and that we are all adults.
No reason needs to be specified for missing a meeting, but the team will flag cases where a member is consistently late or missing meetings.

\subsubsection*{In Case of Emergency}

In the case of an emergency, team members should notify the team as soon as possible. The team will be understanding of emergencies and will work with the member to ensure that they are able to catch up on any missed work.

\subsection*{Accountability and Teamwork}

\subsubsection*{Quality} 

The team will not have any strict standards for the quality of work, but will review each other's work and provide feedback to ensure that any work meets everyone's standards.

\subsubsection*{Attitude}

Team members will approach work with a positive attitude, and any demeaning or negative comments will not be tolerated. Team members will be expected to be respectful and professional in all interactions with each other.

\subsubsection*{Stay on Track}

To keep on track, team members will use Github Projects to track tasks and deliverables.

\subsubsection*{Team Building}

The team will have dinners occassionally.

\subsubsection*{Decision Making} 

Decisions will be made via popular vote.
\end{document}